\documentclass[11pt,a4paper]{article}
\usepackage[T1]{fontenc}
\usepackage[utf8]{inputenc}
\usepackage{lmodern}
\usepackage[margin=26mm,headheight=14pt]{geometry}
\usepackage{amsmath,amssymb,amsthm,mathtools,mathrsfs}
\usepackage{microtype,booktabs,longtable,array,enumitem,needspace}
\usepackage[dvipsnames]{xcolor}
\usepackage{xurl}
\usepackage[colorlinks=true,linkcolor=MidnightBlue,citecolor=MidnightBlue,urlcolor=MidnightBlue]{hyperref}
\usepackage{fancyhdr}
\pagestyle{fancy}\fancyhf{}
\fancyhead[L]{\small Confluent critical transseries}
\fancyhead[R]{\small Charts, local limits, and action budgets}
\fancyfoot[C]{\thepage}
\renewcommand{\headrulewidth}{0.3pt}
\setlength{\parskip}{3pt}
\setlength{\parindent}{1.1em}
\setlength{\emergencystretch}{2em}
\setlist{itemsep=3pt,topsep=5pt}
\setcounter{tocdepth}{2}
\allowdisplaybreaks[2]
\numberwithin{equation}{section}
\newtheorem{theorem}{Theorem}[section]
\newtheorem{proposition}[theorem]{Proposition}
\newtheorem{lemma}[theorem]{Lemma}
\newtheorem{corollary}[theorem]{Corollary}
\theoremstyle{definition}
\newtheorem{definition}[theorem]{Definition}
\newtheorem{example}[theorem]{Example}
\theoremstyle{remark}
\newtheorem{remark}[theorem]{Remark}
% ed. (2026-09-29): unnumbered environment for ProveIt editorial notes; it does not shift any numbering.
\newtheorem*{ednoteinner}{Editorial note (ProveIt, 2026-09-29)}
\newenvironment{ednote}{\begin{ednoteinner}\sloppy\Urlmuskip=0mu plus 1mu\relax}{\end{ednoteinner}}
\newcommand{\R}{\mathbb R}
\newcommand{\C}{\mathbb C}
\newcommand{\N}{\mathbb N}
\newcommand{\Z}{\mathbb Z}
\newcommand{\E}{\mathbb E}
\newcommand{\eps}{\varepsilon}
\newcommand{\dd}{\,\mathrm d}
\newcommand{\ph}{\varphi}
\newcommand{\cE}{\mathcal E}
\newcommand{\cV}{\mathcal V}
\newcommand{\cR}{\mathcal R}
\DeclareMathOperator{\Li}{Li}
\DeclareMathOperator{\Pois}{Pois}
\DeclareMathOperator{\Var}{Var}
\DeclareMathOperator{\supp}{supp}
\hypersetup{pdftitle={Confluent Critical Transseries Across the Exponent-Two Boundary},pdfauthor={Research draft prepared for Vladimir Reshetnikov},pdfsubject={Convergent inverse charts, triangular local limit theorems, conditional Poisson limits, and sharp action budgets}}

\begin{document}
\begin{center}
{\small\scshape Transseries / singular inversion / countable exponential feedback}\par
\vspace{12pt}
{\LARGE\bfseries Confluent Critical Transseries}\par
\vspace{7pt}
{\Large Across the Exponent-Two Boundary}\par
\vspace{9pt}
{\large A Uniform Inverse Chart, Conditional Poisson Limits,\par and Sharp Action Budgets}\par
\vspace{14pt}
{\normalsize Research draft prepared for Vladimir Reshetnikov}\par
{\small 29 September 2026}
\end{center}
\vspace{6pt}
\begin{abstract}
We study the positive feedback equation
\[
 U_\eps(q)=c_\eps\sum_{j\ge1}j^{-3+\eps}q^j e^{jU_\eps(q)},
 \qquad c_\eps=\zeta(2-\eps)^{-1},
\]
in the joint limit in which the coefficient order tends to infinity and
\(\eps\to0\), from either side and at an arbitrary rate. A cancellation-free
normal form yields a single convergent three-variable inverse chart over an
exact implicit core. At \(\eps=0\) this core is a Lambert--\(W_{-1}\) core;
its two outer limits recover fractional-power and finite-variance inversion.
We prove a triangular local Gaussian theorem, stable under arbitrary thinning
above the extreme-action scale. It gives a multiplicative equivalent for the
sector coefficients and a conditional Poisson point-process limit even under
exact lattice conditioning. The fluctuation scale \(B_n\) and extreme scale
\(A_n\) satisfy \(A_n/B_n\to0\) for every \(\eps_n\to0\), while the sharp
retained-coefficient fraction at cutoff \(M\sim sA_n\) is
\(\exp(-1/(2s^2))\). Explicit crossover factors describe the window
\(\eps_n\log n=O(1)\). An exact identity between the inverse core and the
probabilistic normalization connects the two constructions. Complete proofs,
finite real inequalities, exact algebraic checks, and reproducible numerical
diagnostics are included.
\end{abstract}

\noindent\textbf{Research status and boundary.}
This article addresses the \emph{upper-endpoint joint-limit part} of Research
Question~4 in the repository's countable-action critical-Hahn article. The
fixed logarithmic endpoint was already treated in subsequent drafts. The
present contribution is the uniform confluent construction and its
arbitrary-rate, locally conditioned limit theory for the specified model.
The classical Gaussian/extreme-value separation at a fixed cubic tail is not
claimed as new. Publication priority has not been established by an exhaustive
literature review. No Lean formalization, outward-rounded implementation, or
resolution of unrestricted transseries realization is claimed.

\clearpage
\begingroup\small\setlength{\parskip}{0pt}\tableofcontents\endgroup
\clearpage

\section{The research target and the main conclusions}
\label{sec:intro}

\subsection{A precise continuation of the repository program}
The inspected \texttt{ProveIt} research index separates its formal Lean library
from a collection of recent, unmerged research packages \cite{repo-index}.
The article \emph{Beyond Finite-Action Folds} treats positive action weights
with exponent \(1+\alpha\), for fixed \(1<\alpha<2\). Its Research Question~4
asks for joint limits in which \(\alpha\) approaches an endpoint while the
sector order grows, specifically warning against substituting into formulas
containing \(\Gamma(-\alpha)\) \cite{repo-stable}.

Two later drafts treat the fixed upper endpoint, where the tail is exactly
inverse cubic outside a finite prefix \cite{prior-endpoint,prior-marginal}.
The first explicitly excludes a simultaneous moving-exponent limit. We set
\(\alpha=2-\eps\), allow either sign of \(\eps\), and solve a normalized,
exact-tail version of that remaining problem. The restriction to an exact
power tail makes it possible to combine an explicit analytic chart with a
self-contained Fourier proof. A finite-prefix robustness result is proved
separately; no general slowly varying triangular array is silently included.

% ed. (2026-09-29): the two drafts are now filed, and an independent second answer was filed with this package (batch 49).
\begin{ednote}
Both drafts are now filed in this series (batch~48): \cite{prior-endpoint}
as \path{Logarithmic_Critical_Endpoint_Lambert_Charts/} and
\cite{prior-marginal} as \path{Marginal_Critical_Transseries_Action_Budgets/}.
Each lists the present joint limit as its own Question~1; this article
answers the chart and leading-order parts of both, for every rate, without
citing them by number, but not the uniform first cutoff correction that the
marginal draft's question also asks for. The stable--Gaussian package \cite{ed:sge}, filed with this one,
answers the same question independently in the window
\(|\eps|\log n\le K\), without having seen this article. The two agree term
by term: the same scales (its \(B^2=nc_\eps F(\log B)\) and \(N\) are
\eqref{eq:scales-intro}), the same window forms
\eqref{eq:B-window}--\eqref{eq:A-window}, the same minimal budget
\eqref{eq:min-budget} at \(\alpha=2\), the same profile (its
\(\exp(-\ell^{-\alpha}/\alpha)\) tends to \(e^{-1/(2s^2)}\)) and the same
conditioned Poisson limit. Inside the window it goes one order further (a
two-term uniform inverse, the first coefficient correction, a uniform first
cutoff correction without finite-prefix moments, a two-term minimal
budget); this article covers every rate \(\eps_n\to0\) at leading order.
\end{ednote}

The existing manuscripts are used to identify a research target, not as
unverified premises. The coefficient representation, normal form, scale
estimates, local limits, and cutoff conclusions below are proved afresh.
The inspected repository revision and the limits of this comparison are
recorded in Appendix~\ref{app:provenance}.

\subsection{What is proved}
There are two natural problems. The first is local inversion near the critical
value of \(q\); the second is the asymptotic size of \([q^n]U_\eps\). They are
controlled by the same elementary function
\begin{equation}
 h_\eps(x)=\int_1^x y^{-1+\eps}\dd y
 =\begin{cases}(x^\eps-1)/\eps,&\eps\ne0,\\ \log x,&\eps=0.
 \end{cases}
 \label{eq:h-intro}
\end{equation}
The analytic result is an \emph{exact convergent chart}, not merely a finite
asymptotic approximation. For the small positive root of
\(\Psi_\eps(t)=\delta\), it has the form
\begin{equation}
 t=r\,\cV\left(r,\eps,\frac1{h_\eps(1/r)}\right),
 \qquad \frac{c_\eps}{2}r^2h_\eps(1/r)=\delta,
 \label{eq:chart-intro}
\end{equation}
where one holomorphic function \(\cV\) works near the origin in three
independent variables. Its coefficients have a finite recursive construction
and a geometric tail bound.

For any sequence \(\eps_n\to0\), the probabilistic scales are
\begin{equation}
 A_n=(nc_{\eps_n})^{1/(2-\eps_n)},
 \qquad B_n^2=nc_{\eps_n}h_{\eps_n}(B_n),
 \label{eq:scales-intro}
\end{equation}
with the large positive solution chosen for \(B_n\). We prove
\(A_n/B_n\to0\), a uniform lattice local limit on scale \(B_n\), and a
conditional Poisson limit on scale \(A_n\). Consequently,
\begin{equation}
 [q^n]U_{\eps_n}(q)\sim
 \frac{\rho_{\eps_n}^{-n}}{\sqrt{2\pi}\,nB_n},
 \qquad
 \frac{[q^n]U_{\eps_n}^{[M_n]}(q)}{[q^n]U_{\eps_n}(q)}
 \longrightarrow e^{-1/(2s^2)}
 \quad\text{if }M_n/A_n\to s>0.
 \label{eq:main-intro}
\end{equation}
The exponential factor \(\rho_{\eps_n}^{-n}\) must be kept at the actual
parameter value. Substituting \(\rho_0\) there is generally invalid even when
\(\eps_n\to0\).

\subsection{Classical background and novelty qualifications}
Lagrange inversion, Poisson coefficient representations, the analytic implicit
function theorem, and Fourier inversion are standard tools. Janson's
Example~18.29 gives a classical fixed cubic-tail instance of Gaussian
fluctuations decoupling from extremes under conditioning \cite{janson}.
Our point is not to rediscover that endpoint mechanism. It is to prove a
single result for a moving exponent, with no restriction on the rate or sign
of its approach, and to connect it to a convergent singular-inversion chart.
We supply the triangular local-limit argument instead of assuming that a
fixed-distribution theorem is uniform in a changing exponent.

The formulas constitute model-specific theorems with proofs. Whether their
precise combination, or any component in an equivalent normalization, has
appeared elsewhere remains a bibliographic question. In particular,
``confluent'' here describes the merging of the fractional-power and
logarithmic descriptions; it is not a new axiom system for transseries.

\section{The normalized model and exact coefficient identities}
\label{sec:model}
Throughout the real asymptotic arguments, \(|\eps|<1/4\),
\(\alpha=2-\eps\), and
\begin{equation}
 c_\eps=\frac1{\zeta(\alpha)},\qquad
 F_\eps(z)=c_\eps\Li_{1+\alpha}(z)
 =\sum_{j\ge1}a_{\eps,j}z^j,\qquad
 a_{\eps,j}=c_\eps j^{-1-\alpha}.
 \label{eq:model}
\end{equation}
Thus \(\sum j a_{\eps,j}=1\). It is useful to include a coupling parameter:
\begin{equation}
 U_{\eps,\beta}(q)=\beta F_\eps(qe^{U_{\eps,\beta}(q)}),
 \qquad \beta>0.
 \label{eq:feedback}
\end{equation}
The notation without \(\beta\) means \(\beta=1\). A cutoff superscript
\([M]\) means that only actions \(1\le j\le M\) are retained. All formal
solutions have zero constant coefficient. Successive coefficient extraction
in \eqref{eq:feedback} proves their uniqueness.

\begin{proposition}[Lagrange and Poisson identities]
\label{prop:lagrange}
For \(n\ge1\), let \(u_n(\eps,\beta)=[q^n]U_{\eps,\beta}(q)\). Then
\begin{equation}
 u_n(\eps,\beta)=\frac1n[z^n]e^{n\beta F_\eps(z)}.
 \label{eq:lagrange}
\end{equation}
Let \(Z_{n,j}\) be independent Poisson variables of means
\(\nu_{n,j}=n\beta a_{\eps,j}\), and put
\(S_n=\sum jZ_{n,j}\). This sum is finite almost surely, and
\begin{equation}
 u_n(\eps,\beta)=\frac{e^{n\beta F_\eps(1)}}n\Pr(S_n=n).
 \label{eq:poisson}
\end{equation}
If \(S_{n,M}=\sum_{j\le M}jZ_{n,j}\) and
\(\Lambda_{n,M}=\sum_{j>M}\nu_{n,j}\), then
\begin{align}
 R_n(M)&:=\frac{u_n^{[M]}(\eps,\beta)}{u_n(\eps,\beta)}
 =e^{-\Lambda_{n,M}}\frac{\Pr(S_{n,M}=n)}{\Pr(S_n=n)} \label{eq:cutoff-identity}\\
 &=\Pr(Z_{n,j}=0\text{ for all }j>M\mid S_n=n).
 \label{eq:conditional-cutoff}
\end{align}
In particular, \(R_n(M)\) is nondecreasing, belongs to \((0,1]\), and equals
one for \(M\ge n\).
\end{proposition}
\begin{proof}
Set \(z=qe^U\). Then \(z=qe^{\beta F_\eps(z)}\), and the usual formal
Lagrange formula gives
\[
 [q^n]\beta F_\eps(z(q))=
 \frac\beta n[z^{n-1}]F_\eps'(z)e^{n\beta F_\eps(z)}.
\]
Differentiating the exponential proves \eqref{eq:lagrange}. The total Poisson
intensity is finite, so only finitely many Poisson events occur almost surely.
The probability generating function of their weighted sum is
\(\exp(n\beta(F_\eps(z)-F_\eps(1)))\), proving \eqref{eq:poisson}.
Independence of retained and deleted actions gives the two ratio identities.
The positivity conclusions also follow directly from
\eqref{eq:lagrange}; actions exceeding \(n\) cannot affect its \(n\)-th
coefficient.
\end{proof}

At critical coupling define
\begin{equation}
 U_{\eps,c}=F_\eps(1),\qquad
 \rho_\eps=e^{-F_\eps(1)}.
 \label{eq:critical}
\end{equation}
For \(z=e^{-t}\), \(q=\rho_\eps e^{-\delta}\), the exact boundary equation is
\begin{align}
 \delta=\Psi_\eps(t)
 &:=t+F_\eps(e^{-t})-F_\eps(1) \label{eq:psi}\\
 &=c_\eps\sum_{j\ge1}j^{-1-\alpha}
       (e^{-jt}-1+jt), \nonumber\\
 U_{\eps,c}-U_\eps(q)&=t-\delta. \label{eq:udisplacement}
\end{align}
For \(t>0\), termwise differentiation gives
\begin{equation}
 \Psi_\eps'(t)=c_\eps\sum_{j\ge1}j^{-\alpha}(1-e^{-jt})>0,
 \qquad
 \Psi_\eps''(t)=c_\eps\sum_{j\ge1}j^{1-\alpha}e^{-jt}>0.
 \label{eq:psi-derivatives}
\end{equation}
Also \(\Psi_\eps'(t)\to0\) as \(t\downarrow0\), by dominated convergence
using \(\sum j^{-\alpha}<\infty\).

\begin{proposition}[The critical radius]
The radius of convergence of \(U_\eps\) is exactly \(\rho_\eps\), and its
real boundary value is \(U_{\eps,c}\).
\end{proposition}
\begin{proof}
The map \(q(z)=ze^{-F_\eps(z)}\) has positive derivative on \((0,1)\), since
\(zF_\eps'(z)<F_\eps'(1)=1\). It maps that interval onto
\((0,\rho_\eps)\), giving the claimed branch and boundary value. Identity
\eqref{eq:poisson} bounds \(u_n(\eps)\le\rho_\eps^{-n}/n\), so the radius is
at least \(\rho_\eps\). Equations \eqref{eq:udisplacement} and
\eqref{eq:psi-derivatives} show that the derivative of the inverse branch
becomes unbounded at the boundary, since \(\dd t/\dd\delta=1/\Psi_\eps'(t)\).
The branch is therefore not analytic at \(\rho_\eps\), proving equality.
\end{proof}

\section{A cancellation-free normal form}
\label{sec:normal}
Define the entire crossover factor
\begin{equation}
 \cE(y)=\begin{cases}(e^y-1)/y,&y\ne0,\\1,&y=0.\end{cases}
 \qquad h_\eps(x)=\log x\,\cE(\eps\log x).
 \label{eq:E}
\end{equation}
For \(x>1\), let
\begin{equation}
 v_\eps(x)=\frac1{h_\eps(x)},\qquad
 d_\eps(x)=\frac{x^\eps}{h_\eps(x)}=\eps+v_\eps(x).
 \label{eq:vd}
\end{equation}
Both quantities are positive, including when \(\eps<0\).

\begin{lemma}[Uniform elementary estimates]
\label{lem:elementary}
For real \(|\eps|<1/4\) and \(x>1\),
\begin{equation}
 v_\eps(x)+d_\eps(x)\le |\eps|+\frac2{\log x}.
 \label{eq:vd-bound}
\end{equation}
If \(J=\lfloor x\rfloor\), then
\begin{equation}
 0\le\sum_{j=1}^Jj^{-1+\eps}-h_\eps(x)\le1.
 \label{eq:sum-integral}
\end{equation}
Uniformly as \(\eps\to0\) and \(t\downarrow0\),
\begin{align}
 \Psi_\eps(t)&=\frac{c_\eps}{2}t^2h_\eps(1/t)
 \left(1+O\left(|\eps|+\frac1{\log(1/t)}\right)\right), \label{eq:psi-leading}\\
 \Psi_\eps'(t)&=c_\eps t h_\eps(1/t)
 \left(1+O\left(|\eps|+\frac1{\log(1/t)}\right)\right).
 \label{eq:psi-prime-leading}
\end{align}
\end{lemma}
\begin{proof}
Write \(L=\log x\). For \(\eps\ge0\), \(h_\eps(x)\ge L\), so
\(v\le1/L\) and \(d=\eps+v\). For \(\eps<0\), put \(a=-\eps\); then
\(d=a/(e^{aL}-1)\le1/L\) and \(v=a+d\). This proves
\eqref{eq:vd-bound}. The function \(y^{-1+\eps}\) is decreasing. Integrating
on the intervals on either side of each integer gives
\(\int_1^{J+1}f\le\sum_{j\le J}f(j)\le1+\int_1^Jf\), which implies
\eqref{eq:sum-integral} even when \(x\) is not integral.

For the remaining estimates split each series at \(J=\lfloor1/t\rfloor\).
On the lower part,
\(e^{-u}-1+u=u^2/2+O(u^3)\), with an absolute constant for \(0\le u\le1\).
The error is
\(O(t^3\sum_{j\le J}j^{2-\alpha})=O(t^\alpha)\), uniformly in the indicated
exponent interval. On the upper part, \(0\le e^{-u}-1+u\le u\), so its
contribution is also \(O(t^\alpha)\). The quadratic lower sum is
\(t^2(h_\eps(1/t)+O(1))/2\). Relative to the main term the errors are
\(O(v_\eps(1/t)+d_\eps(1/t))\). The derivative follows from
\(1-e^{-u}=u+O(u^2)\) below the split and \(1-e^{-u}\le1\) above it, with
errors \(O(t^{\alpha-1})\). Apply \eqref{eq:vd-bound}.
\end{proof}

The elementary estimates already give the leading crossover. To obtain a
convergent chart we need an exact normal form. Define functions with their
removable values at zero:
\begin{align}
 g(\eps)&=\Gamma(\eps-2)-\frac1{2\eps},&
 b(\eps)&=\zeta(1-\eps)+\frac1\eps, \label{eq:regularparts}\\
 \cR(\eps,t)&=\sum_{k\ge3}
 \frac{(-1)^k\zeta(3-\eps-k)}{k!}\,t^{k-3}. \label{eq:R}
\end{align}
The first values, with \(\gamma_1\) the first Stieltjes constant, are
\begin{align}
 g_0:=g(0)&=\frac34-\frac\gamma2,&
 g_1:=g'(0)&=\frac78-\frac{3\gamma}4+\frac{\gamma^2}4+\frac{\pi^2}{24},
 \label{eq:g-values}\\
 b(0)&=\gamma,& b'(0)&=\gamma_1,& \cR(0,0)&=\frac1{12}. \nonumber
\end{align}
Our sign convention is
\(\zeta(1+s)=1/s+\sum_{k\ge0}(-1)^k\gamma_k s^k/k!\).

\begin{proposition}[Exact confluent normal form]
\label{prop:normal}
The functions \(g\) and \(b\) are holomorphic near zero, and \(\cR\) is
holomorphic on a product neighborhood of \((0,0)\). For small positive \(t\)
and real \(\eps\) near zero,
\begin{equation}
 \frac{\Psi_\eps(t)}{c_\eps}
 =\frac{t^2}{2}\left[h_\eps(1/t)+2g(\eps)t^{-\eps}+b(\eps)\right]
       +t^3\cR(\eps,t).
 \label{eq:exact-normal}
\end{equation}
In particular,
\begin{equation}
 \frac{\Psi_0(t)}{c_0}
 =\frac{t^2}{2}\left(\log\frac1t+\frac32\right)
   +\frac{t^3}{12}-\frac{t^4}{288}+O(t^6),
 \qquad c_0=\frac6{\pi^2}.
 \label{eq:normal-zero}
\end{equation}
\end{proposition}
\begin{proof}
For nonzero \(\eps\) near zero, the standard polylogarithm expansion
\cite[Eq.~25.12.12]{dlmf} gives
\[
 \Li_{3-\eps}(e^{-t})
 =\Gamma(\eps-2)t^{2-\eps}
  +\sum_{k\ge0}\frac{\zeta(3-\eps-k)(-t)^k}{k!},
 \qquad |t|<2\pi,
\]
on the chosen logarithm branch. The constant and linear terms cancel in
\eqref{eq:psi}. Separate the two poles as in \eqref{eq:regularparts}; their
combined contribution is exactly \(t^2h_\eps(1/t)/2\). This proves
\eqref{eq:exact-normal} away from zero. The Laurent expansions of \(\Gamma\)
and \(\zeta\) give \eqref{eq:g-values} and the removable extensions.

For completeness, the tail in \eqref{eq:R} converges locally uniformly in
both variables. The functional equation of \(\zeta\), applied at
\(3-\eps-k\), bounds its coefficient divided by \(k!\) by
\(C k^C(2\pi)^{-k}\) on a sufficiently small closed complex \(\eps\)-disc;
the finitely many initial terms cause no difficulty. The Weierstrass theorem
then gives joint holomorphy for \(|t|<2\pi\) on smaller product domains.
Continuity in \(\eps\) extends \eqref{eq:exact-normal} to zero. Finally,
\(\zeta(0)=-1/2\), \(\zeta(-1)=-1/12\), and \(\zeta(-2)=0\) give
\eqref{eq:normal-zero}.
\end{proof}

\begin{remark}
The separate terms \(\Gamma(\eps-2)t^{2-\eps}\) and
\(\zeta(1-\eps)t^2/2\) have opposite divergent poles. Retaining only the
first one is not uniform when \(\eps\log(1/t)\) remains bounded. The
function \(h_\eps\) retains their cancellation exactly, including at
\(\eps=0\).
\end{remark}

\section{A single convergent inverse chart}
\label{sec:chart}

\subsection{The exact core and its three small parameters}
For \(\delta>0\) small, let \(r=r_\eps(\delta)\in(0,e^{-2})\) solve
\begin{equation}
 \frac{c_\eps}{2}r^2h_\eps(1/r)=\delta,
 \qquad v=\frac1{h_\eps(1/r)}.
 \label{eq:core}
\end{equation}
Existence and uniqueness hold uniformly for \(|\eps|\) sufficiently small.
Indeed, the left side tends to zero at \(r=0\) and its logarithmic derivative
with respect to \(r\) is \(2-d_\eps(1/r)>0\) on this interval, by
\eqref{eq:vd-bound}. At the endpoint,
\begin{equation}
 r_0(\delta)=\exp\left[\frac12W_{-1}\left(-\frac{4\delta}{c_0}\right)\right].
 \label{eq:lambert-core}
\end{equation}
The three variables \((r,\eps,v)\) tend to zero in every joint limit
\(\delta\downarrow0\), \(\eps\to0\), however slowly the exponent moves.

Let
\begin{equation}
 Q_\eps(V)=\frac{V^{-\eps}-1}{\eps},\qquad Q_0(V)=-\log V,
 \label{eq:Q}
\end{equation}
where \(\log V\) is the analytic branch near \(V=1\). This is holomorphic
in \((\eps,V)\) near \((0,1)\).

\begin{theorem}[Convergent confluent chart]
\label{thm:chart}
There is a holomorphic function \(\cV(r,\eps,v)\), on one neighborhood of
\((0,0,0)\), with \(\cV(0,0,0)=1\), such that the unique small positive
solution of \(\Psi_\eps(t)=\delta\) is exactly
\begin{equation}
 t=r\cV(r,\eps,v),\qquad
 r=r_\eps(\delta),\quad v=h_\eps(1/r)^{-1}.
 \label{eq:chart}
\end{equation}
The defining analytic equation for \(V=\cV(r,\eps,v)\) is
\begin{align}
 0={}&V^2\Bigl[1+(\eps+v)Q_\eps(V)
     +2g(\eps)(\eps+v)V^{-\eps}+b(\eps)v \nonumber\\
    &\hspace{38mm}+2rvV\cR(\eps,rV)\Bigr]-1.
 \label{eq:chart-equation}
\end{align}
Its Taylor series is convergent and has a finite coefficient recurrence.
For suitable positive radii \(R_1,R_2,R_3\) and \(K<\infty\), its
truncation \(\cV_{\le N}\) by total degree satisfies
\begin{equation}
 |\cV-\cV_{\le N}|
 \le K\frac{(N+3)^2\theta^{N+1}}{(1-\theta)^3}
 \quad\text{when}\quad
 |r|\le\theta R_1,\quad |\eps|\le\theta R_2,\quad |v|\le\theta R_3,
 \label{eq:chart-tail}
\end{equation}
for \(0<\theta<1\).
\end{theorem}
\begin{proof}
Substitute \(t=rV\) in \eqref{eq:exact-normal} and divide by
\(\delta=c_\eps r^2h_\eps(1/r)/2\). The exact identity
\[
 \frac{h_\eps(1/(rV))}{h_\eps(1/r)}
 =1+(\eps+v)Q_\eps(V),\qquad
 \frac{r^{-\eps}}{h_\eps(1/r)}=\eps+v
\]
gives \eqref{eq:chart-equation}. Its left side, denoted by
\(\mathscr F(r,\eps,v;V)\), is holomorphic in a neighborhood of
\((0,0,0;1)\). At \((r,\eps,v)=(0,0,0)\) it is \(V^2-1\), with derivative
\(2\) at \(V=1\). The analytic implicit function theorem proves the existence
and uniqueness of \(\cV\). For real small parameters this branch is real and
close to one; the positivity and strict monotonicity of \(\Psi_\eps\) identify
it with the desired real root.

Choose a closed polydisc strictly inside the domain of holomorphy, with radii
\(R_i\), on which \(|\cV|\le K\). Cauchy's coefficient estimate bounds the
sum of terms of total degree \(k\) at the stated point by
\(K\binom{k+2}{2}\theta^k\). Summing for \(k>N\), and using
\[
 \sum_{k>N}\binom{k+2}{2}\theta^k
 \le\frac{(N+3)^2\theta^{N+1}}{(1-\theta)^3},
\]
proves \eqref{eq:chart-tail}.

To make the coefficient construction explicit, write
\(\cV=1+\sum_{m\ge1}P_m\), where \(P_m\) is homogeneous of degree \(m\).
If \(P_1,\ldots,P_{m-1}\) are known, then
\begin{equation}
 P_m=-\frac12[s^m]\mathscr F\left(sr,s\eps,sv;
                  1+\sum_{k<m}s^kP_k\right).
 \label{eq:chart-recursion}
\end{equation}
Only finitely many Taylor coefficients enter this extraction. The factor
\(1/2\) comes from the same nonzero derivative, so induction proves the
recurrence at every degree.
\end{proof}

\subsection{The first corrections and their support}
The first two homogeneous terms are
\begin{align}
 P_1={}&-g_0\eps-\frac34v, \label{eq:P1}\\
 P_2={}&\left(\frac32g_0^2-\frac12g_0-g_1\right)\eps^2
 +\left(\frac74g_0-\frac38-g_1-\frac12\gamma_1\right)\eps v \nonumber\\
 &+\frac{15}{32}v^2-\frac1{12}rv.
 \label{eq:P2}
\end{align}
For example, the endpoint specialization is
\begin{equation}
 \frac{t}{r}=1-\frac34v+\frac{15}{32}v^2-\frac1{12}rv
       +O\bigl((|r|+|v|)^3\bigr),\qquad \eps=0.
 \label{eq:endpoint-chart}
\end{equation}
The formulas follow by inserting \(V=1+P_1+P_2\) into
\eqref{eq:chart-equation} and collecting total degrees. The supplied exact
symbolic check verifies that the residual through degree two is identically
zero, treating the constants as symbols subject only to
\(2g_0+\gamma=3/2\).

The support before evaluation is the ordinary locally finite semigroup
\(\N^3\): terms are \(r^i\eps^jv^k\). After evaluation, \(v\) is a
nontrivial logarithmic/exponential function of \(r\) and \(\eps\). Thus the
theorem gives one convergent analytic chart \emph{over an exact core}; it does
not assert that the resulting family has one finite Puiseux or Hahn grid in
\(\delta\), independent of the parameter. That distinction is what permits
confluence without introducing a false uniform ramification assertion.

The constants in \eqref{eq:chart-tail} exist by holomorphy; numerical values
for a certified polydisc have not been computed here. The separate real
inequalities in Section~\ref{sec:certificates} provide finite root brackets
without requiring those constants.

\subsection{Recovery of the two fixed-exponent boundary laws}
The same chart also displays the two limits when a sufficiently small nonzero
\(\eps\) is held fixed while \(r\downarrow0\). For \(\eps>0\), one has
\(v\to0\), and \eqref{eq:chart-equation} at \((r,v)=(0,0)\) reduces to
\[
 V^{2-\eps}\bigl(1+2\eps g(\eps)\bigr)=1.
\]
For \(\eps<0\), one has \(v\to-\eps\), and the limiting equation is
\(V^2(1-\eps b(\eps))=1\). Hence, by the branch near one,
\begin{align}
 \cV(0,\eps,0)&=[2\eps\Gamma(\eps-2)]^{-1/(2-\eps)}
       &&(\eps>0), \label{eq:stable-chart-limit}\\
 \cV(0,\eps,-\eps)&=[-\eps\zeta(1-\eps)]^{-1/2}
       &&(\eps<0). \label{eq:gauss-chart-limit}
\end{align}
Combining these identities with the core gives the familiar leading forms
\begin{equation}
 t\sim
 \begin{cases}
 [\delta/(c_\eps\Gamma(\eps-2))]^{1/(2-\eps)},&\eps>0,\\[2pt]
 [2\delta/(c_\eps\zeta(1-\eps))]^{1/2},&\eps<0.
 \end{cases}
 \label{eq:fixed-exponents}
\end{equation}
The negative-\(\eps\) side has a finite second moment and a leading square
root. It is not an analytic finite-action fold: higher moments and higher
boundary derivatives can still be singular.

\section{The two scales in an arbitrary joint limit}
\label{sec:scales}
Now let \(\eps_n\to0\), with either sign and without a rate restriction. Write
\(\alpha_n=2-\eps_n\), \(c_n=c_{\eps_n}\), and define \(A_n,B_n\) by
\eqref{eq:scales-intro}. All statements below concern sufficiently large
\(n\), so the large root \(B_n>e^2\) is unambiguous. Set
\begin{equation}
 v_n=\frac1{h_{\eps_n}(B_n)},\qquad
 d_n=\frac{B_n^{\eps_n}}{h_{\eps_n}(B_n)}
     =nc_nB_n^{-\alpha_n}=\eps_n+v_n.
 \label{eq:dn}
\end{equation}

\begin{lemma}[Scale separation and retained truncated variance]
\label{lem:scales}
The scales satisfy
\begin{equation}
 A_n=n^{1/2+o(1)},\qquad B_n=n^{1/2+o(1)},\qquad
 v_n,d_n\to0,\qquad \frac{A_n}{B_n}=d_n^{1/\alpha_n}\to0.
 \label{eq:scale-separation}
\end{equation}
For every fixed \(a>0\),
\begin{equation}
 \frac{h_{\eps_n}(aA_n)}{h_{\eps_n}(B_n)}\longrightarrow1,
 \qquad
 nc_n\sum_{j>aA_n}j^{-\alpha_n}=O_a(A_n)=o(B_n).
 \label{eq:truncated-variance}
\end{equation}
The first assertion is uniform for \(a\) in a compact subset of \((0,\infty)\).
\end{lemma}
\begin{proof}
The function \(x^2/h_\eps(x)\) is increasing on \(x>e^2\) for small
\(|\eps|\), and tends to infinity, which proves uniqueness and existence of
\(B_n\) there. It also shows \(B_n\to\infty\). Bound
\eqref{eq:vd-bound} then proves \(v_n,d_n\to0\). Moreover
\[
 h_\eps(x)\le(\log x)e^{\eps_+\log x},\qquad \eps_+=\max(\eps,0),
\]
so \(\log h_{\eps_n}(B_n)=o(\log B_n)\); here \(h\to\infty\), so no
negative lower-order logarithm is needed. Taking logarithms in the defining
equation gives \(\log B_n/\log n\to1/2\). The assertion for \(A_n\) is
direct, and \((A_n/B_n)^{\alpha_n}=d_n\) is exact.

For the truncated-variance assertion, put \(B=B_n\), \(A=A_n\),
\(\eps=\eps_n\), \(d=d_n\), and
\(\Delta=\log(B/(aA))\); eventually \(aA<B\). If \(\eps\ge0\),
\[
 \frac{h_\eps(B)-h_\eps(aA)}{h_\eps(B)}
 =\frac1{h_\eps(B)}\int_{aA}^B y^\eps\frac{\dd y}{y}
 \le d\Delta.
\]
If \(\eps<0\), the same integral is bounded by
\[
 \frac{(aA)^\eps}{h_\eps(B)}\Delta
 =a^\eps d^{\,1+\eps/\alpha_n}\Delta
 =a^\eps d^{\,2/\alpha_n}\Delta.
\]
Since \(\Delta=-\log a-(\log d)/\alpha_n\), both bounds tend to zero,
uniformly for \(a\) in compact positive intervals. Finally, the integral
bound for the decreasing function \(y^{-\alpha_n}\), together with
\(nc_n=A_n^{\alpha_n}\), bounds the deleted mean by a constant times
\(nc_n(aA_n)^{1-\alpha_n}=O_a(A_n)\).
\end{proof}

The second-moment sum up to \(B_n\) is of order \(h_{\eps_n}(B_n)\).
Lemma~\ref{lem:scales} says that retaining only actions up to a fixed multiple
of \(A_n\), though \(A_n\ll B_n\), still retains asymptotically all of this
quantity. This fact is stronger than scale separation alone and is the key
to the local conditioning theorem.

\begin{theorem}[Explicit scale crossover]
\label{thm:scale-crossover}
Let \(\lambda_n=\eps_n\log n\). If \(|\lambda_n|\le K\), then uniformly
for such sequences,
\begin{align}
 B_n^2&=\frac{c_n}{2}n\log n\,\cE(\lambda_n/2)
       \left(1+O_K\left(\frac{\log\log n}{\log n}\right)\right),
       \label{eq:B-window}\\
 A_n&=\sqrt{nc_n}\,e^{\lambda_n/4}
       \left(1+O_K\left(\frac1{\log n}\right)\right).
       \label{eq:A-window}
\end{align}
In particular, if \(\lambda_n\to\lambda\in\R\),
\begin{equation}
 \frac{B_n}{A_n}\sim
 \sqrt{\log n\,\frac{1-e^{-\lambda/2}}{\lambda}},
 \label{eq:separation-window}
\end{equation}
where the fraction at zero equals \(1/2\). In the outer joint regimes,
\begin{align}
 B_n&\sim(nc_n/\eps_n)^{1/(2-\eps_n)}
 &&\text{if }\eps_n\log n\to+\infty, \label{eq:B-positive}\\
 B_n&\sim\sqrt{nc_n/|\eps_n|}
 &&\text{if }\eps_n\log n\to-\infty. \label{eq:B-negative}
\end{align}
\end{theorem}
\begin{proof}
When \(|\lambda_n|\le K\), the defining equation and
\(h_\eps(x)=\log x\,\cE(\eps\log x)\) first give
\(\log B_n=\tfrac12\log n+O_K(\log\log n)\). On compact real intervals
\(\cE\) is positive with bounded logarithmic derivative. Substitution gives
\eqref{eq:B-window}. Expanding \(1/(2-\eps_n)\) in the exact expression for
\(A_n\) gives \eqref{eq:A-window}. Their quotient gives
\eqref{eq:separation-window}.

In the positive outer regime, \(\eps_n\log B_n\to+\infty\), by
Lemma~\ref{lem:scales}. Consequently
\(h_{\eps_n}(B_n)\sim B_n^{\eps_n}/\eps_n\). Raising the resulting
relation for \(B_n^{2-\eps_n}\) to the bounded power \(1/(2-\eps_n)\)
proves \eqref{eq:B-positive}. In the negative regime,
\(h_{\eps_n}(B_n)\sim1/|\eps_n|\), proving \eqref{eq:B-negative}.
\end{proof}

\section{A local Gaussian theorem stable under tail thinning}
\label{sec:llt}
Let \(\ph(x)=(2\pi)^{-1/2}e^{-x^2/2}\). In this section the unthinned
Poisson means are \(nc_nj^{-1-\alpha_n}\), so \(\E S_n=n\).

\begin{theorem}[Masked triangular local limit]
\label{thm:masked-llt}
Fix \(a>0\). For each \(n\), choose arbitrary numbers
\(0\le w_{n,j}\le1\) satisfying \(w_{n,j}=1\) for all \(j\le aA_n\).
Let \(T_n\) be the weighted sum of independent Poisson variables with means
\(w_{n,j}nc_nj^{-1-\alpha_n}\), and let \(\mu_n=\E T_n\). Then
\begin{equation}
 \sup_{k\in\Z}\left|
 B_n\Pr(T_n=k)-\ph\left(\frac{k-\mu_n}{B_n}\right)
 \right|\longrightarrow0.
 \label{eq:masked-llt}
\end{equation}
The convergence is uniform over the indicated masks, and
\(\mu_n=n+O_a(A_n)=n+o(B_n)\). In particular,
\begin{equation}
 \sup_{k\in\Z}\left|
 B_n\Pr(S_n=k)-\ph\left(\frac{k-n}{B_n}\right)
 \right|\longrightarrow0.
 \label{eq:llt}
\end{equation}
\end{theorem}

\subsection{Convergence on compact Fourier intervals}
\begin{proof}[Proof of Theorem~\ref{thm:masked-llt}, first part]
Suppress the subscript \(n\), and write \(B=B_n\), \(A=A_n\),
\(\eps=\eps_n\), \(\alpha=\alpha_n\), \(c=c_n\), \(h=h_\eps(B)\).
The centered unthinned logarithmic characteristic function at frequency
\(x/B\) is
\begin{equation}
 K_n(x)=nc\sum_{j\ge1}j^{-1-\alpha}
        \left(e^{ijx/B}-1-\frac{ijx}{B}\right).
 \label{eq:K}
\end{equation}
On every fixed interval \(|x|\le T\), Taylor's formula for \(j\le B\) gives
\begin{equation}
 K_n(x)=-\frac{x^2}{2}\frac{\sum_{j\le B}j^{1-\alpha}}h+O_T(d_n)
       =-\frac{x^2}{2}+O_T(v_n+d_n).
 \label{eq:K-compact}
\end{equation}
Indeed the lower-part cubic error is bounded by
\[
 C_T\frac{nc}{B^3}\sum_{j\le B}j^{2-\alpha}
 \le C_TncB^{-\alpha}=C_Td_n.
\]
For \(j>B\), the bound
\(\lvert e^{iu}-1-iu\rvert\le2+|u|\) gives the same order. The sum-integral
estimate \eqref{eq:sum-integral} accounts for \(v_n\).

The masked centered logarithm has an extra factor \(w_{n,j}\) in the sum.
For deleted terms with \(aA<j\le B\), the inequality
\(\lvert e^{iu}-1-iu\rvert\le u^2/2\) bounds the difference by
\[
 C_T\frac{\sum_{aA<j\le B}j^{1-\alpha}}h.
\]
For \(j>B\), it is \(O_T(d_n)\). Lemma~\ref{lem:scales} and
\eqref{eq:sum-integral} show that the difference tends to zero, independently
of the mask. Thus the centered masked characteristic functions converge to
\(e^{-x^2/2}\), uniformly on compact \(x\)-intervals. The mean estimate is
exactly the deleted-mean bound in \eqref{eq:truncated-variance}.
\end{proof}

\subsection{A uniform integrable Fourier majorant}
Compact convergence is not enough for a local limit. We next bound the full
lattice inversion integral, including frequencies that grow with \(n\).

\begin{lemma}[Two Fourier bounds]
\label{lem:fourier}
There are positive constants \(C_1,C_2\), depending at most on \(a\), such
that all sufficiently large \(n\), and every mask in
Theorem~\ref{thm:masked-llt}, satisfy
\begin{align}
 \left|\E e^{ix(T_n-\mu_n)/B_n}\right|
 &\le e^{-C_1|x|^{7/4}}
 &&(1\le|x|\le\sqrt{B_n}), \label{eq:fourier-near}\\
 \left|\E e^{ix(T_n-\mu_n)/B_n}\right|
 &\le e^{-C_2n/B_n}
 &&(\sqrt{B_n}\le|x|\le\pi B_n).
 \label{eq:fourier-far}
\end{align}
\end{lemma}
\begin{proof}
The modulus is
\[
 \exp\left[-nc\sum_{j\ge1}w_{n,j}j^{-1-\alpha}
                   (1-\cos(jx/B))\right].
\]
For the first range put \(y=\min(aA,B/|x|)\). All integers \(j\le y\)
are retained and satisfy \(|jx/B|\le1\). Since
\(1-\cos u\ge u^2/4\) there,
\begin{equation}
 -\log\left|\E e^{ix(T_n-\mu_n)/B}\right|
 \ge\frac{x^2}{4h}\sum_{j\le y}j^{1-\alpha}
 \ge\frac{x^2}{4}\frac{h_\eps(y)}h.
 \label{eq:fourier-basic}
\end{equation}
If \(y=aA\), the ratio tends to one by Lemma~\ref{lem:scales}. If
\(y=B/|x|\), write \(L=\log B\), \(b=\log|x|\le L/2\).
For \(\eps\ge0\), changing variables in the defining integral shows
\[
 h_\eps(B/|x|)
 =e^{-\eps b}\int_b^L e^{\eps u}\dd u
 \ge\tfrac12 e^{-\eps b}h_\eps(B).
\]
The last inequality follows because the increasing integrand has at least
half its integral on the upper interval of length at least \(L/2\).
For \(\eps<0\), the decreasing integrand gives
\(\int_0^{L-b}e^{\eps u}\dd u\ge\tfrac12\int_0^Le^{\eps u}\dd u\).
Therefore in either case
\[
 \frac{h_\eps(B/|x|)}{h_\eps(B)}
 \ge\frac12|x|^{-\eps_+}\ge\frac12|x|^{-1/4}.
\]
Together with \eqref{eq:fourier-basic} this proves
\eqref{eq:fourier-near}, with a smaller constant if necessary.

For the second range the retained action \(j=1\) alone suffices. On
\(|u|\le\pi\), \(1-\cos u\ge2u^2/\pi^2\). The exponent is consequently
at least \(C n x^2/B^2\), which is at least \(C n/B\) on the stated range.
The constants are uniform because \(c_n\to c_0>0\).
\end{proof}

\begin{proof}[Completion of Theorem~\ref{thm:masked-llt}]
Extend the centered scaled characteristic function by zero outside
\([ -\pi B_n,\pi B_n]\). Compact convergence, the integrable majorant in
\eqref{eq:fourier-near}, and
\[
 B_ne^{-C_2n/B_n}\longrightarrow0
\]
from \(B_n=n^{1/2+o(1)}\), prove convergence in \(L^1(\R)\) to
\(e^{-x^2/2}\). Every bound is independent of the mask. Lattice Fourier
inversion, after the change of variables \(x=B_n\theta\), gives
\[
 B_n\Pr(T_n=k)=\frac1{2\pi}\int_{-\pi B_n}^{\pi B_n}
 e^{-ix(k-\mu_n)/B_n}\E e^{ix(T_n-\mu_n)/B_n}\dd x.
\]
The difference from the Gaussian inversion integral is bounded uniformly
in \(k\) by the \(L^1\) difference divided by \(2\pi\). This proves the
claim and completes the proof without an unproved interchange of limits.
\end{proof}

\begin{remark}[Meaning of arbitrary-rate uniformity]
The theorem applies to every sequence \(\eps_n\to0\), including sign changes
and \(|\eps_n|\log n\to\infty\). Equivalently, its errors tend to zero
uniformly over any prescribed shrinking neighborhoods
\(|\eps|\le\eta_n\to0\). It does not assert Gaussian convergence uniformly
on a fixed exponent neighborhood: a fixed \(\eps>0\) belongs to a
non-Gaussian stable regime.
\end{remark}

\section{Sector asymptotics and sharp action budgets}
\label{sec:budgets}

\begin{theorem}[Multiplicative sector equivalent]
\label{thm:sectors}
For every \(\eps_n\to0\),
\begin{equation}
 u_n(\eps_n)=\frac{\rho_{\eps_n}^{-n}}{\sqrt{2\pi}\,nB_n}(1+o(1)).
 \label{eq:sectors}
\end{equation}
If \(\eps_n\log n\to\lambda\in\R\), then
\begin{equation}
 u_n(\eps_n)\sim
 \frac{\rho_{\eps_n}^{-n}}
 {\sqrt{\pi c_0}\,n^{3/2}\sqrt{\log n}\sqrt{\cE(\lambda/2)}}.
 \label{eq:sectors-window}
\end{equation}
\end{theorem}
\begin{proof}
At critical coupling the mean of \(S_n\) is exactly \(n\). Theorem
\ref{thm:masked-llt} gives \(B_n\Pr(S_n=n)\to\ph(0)>0\), so its additive
local limit is a relative equivalent at that point. Substitute into
\eqref{eq:poisson}, and then use \eqref{eq:B-window}.
\end{proof}

\begin{theorem}[Cutoff profile and necessary-and-sufficient budget]
\label{thm:budget}
Define
\begin{equation}
 H(s)=e^{-1/(2s^2)}\quad(s>0),\qquad H(0)=0,\quad H(\infty)=1.
 \label{eq:H-profile}
\end{equation}
For any positive integer sequence \(M_n\), if \(M_n/A_n\to s\) in
\([0,\infty]\), then
\begin{equation}
 R_n(M_n)\longrightarrow H(s).
 \label{eq:budget-profile}
\end{equation}
The convergence is uniform when \(M_n/A_n\) ranges over a compact positive
interval, with the profile evaluated at the actual ratio. Moreover,
\begin{equation}
 R_n(M_n)\to1\quad\Longleftrightarrow\quad M_n/A_n\to\infty.
 \label{eq:iff-budget}
\end{equation}
For fixed allowed loss \(\theta\in(0,1)\), let
\(M_{\min}(n,\theta)=\min\{M:R_n(M)\ge1-\theta\}\). Then
\begin{equation}
 \frac{M_{\min}(n,\theta)}{A_n}\longrightarrow
 \frac1{\sqrt{2\log(1/(1-\theta))}}.
 \label{eq:min-budget}
\end{equation}
\end{theorem}
\begin{proof}
First suppose \(M_n/A_n\to s\in(0,\infty)\). An integral comparison, uniform
for \(\alpha_n\) near two, gives
\begin{equation}
 \Lambda_{n,M_n}=nc_n\sum_{j>M_n}j^{-1-\alpha_n}
 =\frac{nc_n}{\alpha_nM_n^{\alpha_n}}(1+O(M_n^{-1}))
 \longrightarrow\frac1{2s^2}.
 \label{eq:lambda-limit}
\end{equation}
The cutoff mask retains every action up to \(aA_n\), for any fixed
\(0<a<s\) and all sufficiently large \(n\). Its mean differs from \(n\) by
\(O(A_n)=o(B_n)\). The masked local limit thus gives
\(B_n\Pr(S_{n,M_n}=n)\to\ph(0)\), as for the full sum. The exact ratio
\eqref{eq:cutoff-identity} proves \eqref{eq:budget-profile}.

Uniformity on compact positive intervals follows either from the uniform
mask bounds or from the usual monotonicity argument: choose a finite mesh,
apply pointwise convergence on that mesh, and use uniform continuity of the
continuous increasing function \(H\) between adjacent mesh points.

If \(M_n/A_n\to0\), compare \(M_n\) to \(\lfloor sA_n\rfloor\) for every
fixed \(s>0\), and then let \(s\downarrow0\). If \(M_n/A_n\to\infty\),
compare in the opposite direction and let \(s\to\infty\). This proves the
endpoint statements. Conversely, if the ratio of cutoffs does not tend to
infinity, some subsequence is bounded above by a fixed \(C\). On that
subsequence \(R_n(M_n)\le R_n(\lceil CA_n\rceil)\to H(C)<1\), proving
necessity in \eqref{eq:iff-budget}.

Finally, the strictly increasing continuous profile crosses \(1-\theta\)
only at the value displayed in \eqref{eq:min-budget}. Compare cutoffs just
below and just above that value times \(A_n\), and use monotonicity and
\(A_n\to\infty\) to remove integer rounding.
\end{proof}

\begin{corollary}[The explicit exponent/order window]
\label{cor:window-budget}
If \(\eps_n\log n\to\lambda\in\R\) and
\(M_n/\sqrt{nc_0}\to m>0\), then
\begin{equation}
 R_n(M_n)\longrightarrow
 \exp\left(-\frac{e^{\lambda/2}}{2m^2}\right).
 \label{eq:window-budget}
\end{equation}
Thus a budget appropriate to the endpoint \(\lambda=0\) does not retain the
same fraction throughout the crossover window.
\end{corollary}
\begin{proof}
Equation \eqref{eq:A-window} gives
\(M_n/A_n\to me^{-\lambda/4}\). Apply Theorem~\ref{thm:budget}.
\end{proof}

These are statements about a \emph{prefix cutoff}. No optimality among all
possible non-prefix subsets of actions is claimed. Likewise,
\eqref{eq:min-budget} is proved for a fixed loss tolerance, not for a tolerance
\(\theta_n\) tending to zero. The latter requires quantitative errors.

\section{Conditional Poisson limits for the entire extreme-action process}
\label{sec:ppp}
Define the random point measure
\begin{equation}
 \Pi_n=\sum_{j\ge1}Z_{n,j}\,\delta_{j/A_n}
 \quad\text{on }(0,\infty).
 \label{eq:pointprocess}
\end{equation}
It is locally finite. A Poisson point process with intensity
\(y^{-3}\dd y\) is characterized by the Laplace functional
\begin{equation}
 \exp\left[-\int_0^\infty(1-e^{-f(y)})y^{-3}\dd y\right],
 \qquad f\ge0,\quad f\in C_c((0,\infty)).
 \label{eq:ppp-laplace}
\end{equation}
The infinite intensity near zero causes no problem for compactly supported
test functions in \((0,\infty)\).

\begin{theorem}[Local factorization under exact conditioning]
\label{thm:ppp}
Let \(\eps_n\to0\), and let integers \(k_n\) satisfy
\((k_n-n)/B_n\to x\in\R\). For every nonnegative
\(f\in C_c((0,\infty))\),
\begin{align}
 &B_n\E\left[e^{-\Pi_n(f)}\mathbf1_{\{S_n=k_n\}}\right] \nonumber\\
 &\qquad\longrightarrow
 \ph(x)\exp\left[-\int_0^\infty(1-e^{-f(y)})y^{-3}\dd y\right].
 \label{eq:local-factorization}
\end{align}
Consequently, under \(S_n=k_n\), the point process \(\Pi_n\) converges to
that Poisson point process. The limit is independent of the limiting central
location \(x\). Unconditionally, \(((S_n-n)/B_n,\Pi_n)\) converges jointly
to an independent standard normal variable and the same Poisson process.
\end{theorem}
\begin{proof}
Choose \(a>0\) below the support of \(f\), and set
\(w_{n,j}=e^{-f(j/A_n)}\). Poisson exponential tilting, first for finitely
many nontrivial weights and then directly for the present compact support,
gives the exact identity
\begin{equation}
 \E[e^{-\Pi_n(f)}\mathbf1_{\{S_n=k\}}]
 =e^{-L_n(f)}\Pr(T_n=k),\qquad
 L_n(f)=\sum_jnc_nj^{-1-\alpha_n}(1-w_{n,j}),
 \label{eq:poisson-tilt}
\end{equation}
where \(T_n\) uses the masked means from
Theorem~\ref{thm:masked-llt}. The identity follows already for one Poisson
variable from multiplying its mass at \(m\) by \(w^m\); independence gives
the product.

Since \(nc_nA_n^{-\alpha_n}=1\), a Riemann sum on the compact support of
\(f\), bounded away from zero, yields
\begin{equation}
 L_n(f)=\frac1{A_n}\sum_j
 (j/A_n)^{-1-\alpha_n}(1-e^{-f(j/A_n)})
 \longrightarrow\int_0^\infty(1-e^{-f(y)})y^{-3}\dd y.
 \label{eq:tilt-riemann}
\end{equation}
The masked mean is \(n+O(A_n)=n+o(B_n)\). Its local limit therefore gives
\(B_n\Pr(T_n=k_n)\to\ph(x)\). Combine this with
\eqref{eq:poisson-tilt} to obtain \eqref{eq:local-factorization}. Dividing by
\(B_n\Pr(S_n=k_n)\to\ph(x)>0\) proves convergence of the conditional
Laplace functionals, which characterizes the point-process limit.

For the unconditional joint assertion, the same proof is uniform for
\((k-n)/B_n\) in a fixed compact interval. Sum
\eqref{eq:poisson-tilt} against a bounded continuous function of
\((k-n)/B_n\); the local limit gives a Gaussian Riemann integral on that
interval. The unmasked central limit, already implied by compact
characteristic-function convergence, controls the complement. This proves
factorization for products of a central test function and a point-process
Laplace functional, and hence the stated joint convergence.
\end{proof}

The local assertion is stronger than an unconditional independence statement:
the event \(S_n=k_n\) has probability of order \(1/B_n\). An unconditional
weak limit alone would not justify dividing by that vanishing probability.
The masked local limit is precisely what makes that division legitimate.

\begin{corollary}[Ordered extreme actions]
\label{cor:ordered}
Let \(J_{n,(1)}\ge J_{n,(2)}\ge\cdots\) be the action sizes in the Poisson
multiset, including multiplicities and padded by zeros if necessary. Under
\(S_n=n\), for each fixed \(r\ge1\) and \(s>0\),
\begin{equation}
 \Pr\left(J_{n,(r)}/A_n\le s\mid S_n=n\right)
 \longrightarrow e^{-\mu(s)}\sum_{\ell=0}^{r-1}\frac{\mu(s)^\ell}{\ell!},
 \qquad \mu(s)=\frac1{2s^2}.
 \label{eq:ordered}
\end{equation}
\end{corollary}
\begin{proof}
On any bounded interval bounded away from zero, Theorem~\ref{thm:ppp} gives
the corresponding Poisson counting limit; boundary points have zero limiting
intensity. To pass from \((s,R]\) to \((s,\infty)\), observe from
Theorem~\ref{thm:budget} that the conditional probability of any action
exceeding \(RA_n\) tends to \(1-H(R)\), which vanishes as \(R\to\infty\).
The limiting number above \(s\) is Poisson of mean
\(\int_s^\infty y^{-3}\dd y=1/(2s^2)\). The \(r\)-th largest is at most
\(s\) precisely when this count is at most \(r-1\).
\end{proof}

\section{A moving coupling window and finite-prefix robustness}
\label{sec:coupling}

\begin{theorem}[A bounded critical coupling window]
\label{thm:coupling}
Let \(\eps_n\to0\), and let \(\omega_n\) range over a fixed compact subset
of \(\R\). Define
\begin{equation}
 \beta_n=1+\omega_n\frac{B_n}{n}.
 \label{eq:beta-window}
\end{equation}
Then \(\beta_n>0\) for all sufficiently large \(n\), and uniformly over
this window,
\begin{equation}
 u_n(\eps_n,\beta_n)
 =\frac{e^{n\beta_nF_{\eps_n}(1)}}{\sqrt{2\pi}\,nB_n}
    e^{-\omega_n^2/2}(1+o(1)).
 \label{eq:coupling-sector}
\end{equation}
For \(M_n/A_n\to s>0\), the cutoff ratio remains \(H(s)\). Conditioned on
the total weighted action being \(n\), the point-process limit remains the
Poisson process with intensity \(y^{-3}\dd y\).
\end{theorem}
\begin{proof}
The factor \(\beta_n\to1\) multiplies every Poisson mean. The centered
characteristic exponent and each Fourier lower bound in the proof of
Theorem~\ref{thm:masked-llt} are multiplied by this factor. All compact
limits and integrable bounds therefore remain valid uniformly in
\(\omega_n\). The unmasked mean is \(n\beta_n=n+\omega_nB_n\). At the
target lattice point \(n\), the centered coordinate is exactly
\(-\omega_n\), giving the local factor \(\ph(\omega_n)\), bounded away from
zero on the compact window. Equation \eqref{eq:poisson} proves
\eqref{eq:coupling-sector}.

Deletion above a fixed multiple of \(A_n\) still shifts the mean by only
\(O(A_n)=o(B_n)\). Thus numerator and denominator in the cutoff ratio have
the same Gaussian factor. The deleted intensity has the same limit because
\(\beta_n\to1\). Similarly, the tilted Riemann sum and the ratio of local
masses in \eqref{eq:poisson-tilt} are unchanged in the limit, proving the
point-process assertion.
\end{proof}

For \(\beta_n\ne1\), the carrier \(e^{n\beta_nF_{\eps_n}(1)}\) in
\eqref{eq:coupling-sector} is an exact factor from coefficient extraction.
It is not being identified with the reciprocal radius of convergence raised
to the \(n\)-th power. In particular, on the supercritical side the dominant
singularity can move into the analytic interior. The bounded-window proof
does not require a separate assertion about that singularity.

\begin{corollary}[Robustness under a bounded finite prefix]
\label{cor:prefix}
The probabilistic results of Sections~\ref{sec:llt}--\ref{sec:coupling} remain
valid for weights \(a_{n,j}\ge0\) satisfying, for some fixed \(J\),
\begin{equation}
 a_{n,j}=c_nj^{-3+\eps_n}\quad(j>J),\qquad
 \sum_{j\ge1}j a_{n,j}=1,
 \label{eq:prefix-model}
\end{equation}
where \(\eps_n\to0\), \(c_n\to c_*>0\), the finitely many prefix weights
are bounded, and \(a_{n,1}\ge a_*>0\). Define \(A_n,B_n\) using this \(c_n\),
and set \(\rho_n=\exp(-\sum_j a_{n,j})\). In the explicit bounded-window
formulas replace \(c_0\) by \(c_*\).
\end{corollary}
\begin{proof}
The scale lemmas use only \(c_n\to c_*>0\). After exact centering, altering a
fixed bounded prefix changes the compact characteristic exponent by
\(O(n/B_n^2)=O(1/h_{\eps_n}(B_n))=o(1)\). The tail integrals, mean deletion
bounds, and Riemann sums are unchanged. In the near-frequency bound, omitting
finitely many terms changes the diverging second-moment sum by \(O(1)\),
which is negligible even at the smallest truncation used there. The far
bound uses \(a_{n,1}\ge a_*>0\) instead of \(c_n\). The Lagrange and Poisson
identities hold for arbitrary nonnegative weights, completing the proof.
\end{proof}

This corollary concerns the asymptotic coefficient and point-process results.
The explicit chart coefficients \eqref{eq:P1}--\eqref{eq:P2} belong to the
normalized polylogarithm model; finite-prefix changes generally alter them.

\section{Inverse crossover profiles and an exact core duality}
\label{sec:duality}

\begin{theorem}[Real inverse crossover]
\label{thm:inverse-window}
Let \(\delta\downarrow0\), \(\eps\to0\), and let \(t>0\) solve
\(\Psi_\eps(t)=\delta\). Write \(L=\log(1/\delta)\). Uniformly when
\(|\eps L|\le K\),
\begin{equation}
 t=2\sqrt{\frac{\delta}{c_\eps L\cE(\eps L/2)}}
       \left(1+O_K\left(\frac{\log L}{L}\right)\right).
 \label{eq:inverse-window}
\end{equation}
In the outer joint regimes,
\begin{align}
 t&\sim(2\eps\delta/c_\eps)^{1/(2-\eps)}
 &&\text{if }\eps L\to+\infty, \label{eq:inverse-positive}\\
 t&\sim\sqrt{2|\eps|\delta/c_\eps}
 &&\text{if }\eps L\to-\infty. \label{eq:inverse-negative}
\end{align}
At \(\eps=0\), the leading form is
\(t\sim2\sqrt{\delta/(c_0\log(1/\delta))}\). In all these limits,
\(U_{\eps,c}-U_\eps(\rho_\eps e^{-\delta})\sim t\).
\end{theorem}
\begin{proof}
The chart gives \(t/r\to1\) in every joint limit. Put
\(\ell=\log(1/r)\). Taking logarithms of the core equation shows
\(\ell\sim L/2\), by the same estimates used for \(\log B_n\). If
\(|\eps L|\le K\), refine this to
\(\ell=L/2+O_K(\log L)\). Hence
\[
 h_\eps(1/r)=\frac L2\cE(\eps L/2)
              \left(1+O_K\left(\frac{\log L}{L}\right)\right).
\]
The core equation gives the displayed formula for \(r\); the chart error is
\(O(|\eps|+v+r)=O_K(1/L)\), which is smaller than the retained error.

When \(\eps L\to+\infty\), one has
\(h_\eps(1/r)\sim r^{-\eps}/\eps\); on the negative side it is
\(\sim1/|\eps|\). Solve the core equation in each case and use \(t/r\to1\).
Finally, \(\delta/t=\Psi_\eps(t)/t\to0\) uniformly in this joint limit:
\eqref{eq:psi-leading} and
\(h_\eps(1/t)\le\log(1/t)t^{-\eps_+}\) make this immediate for
\(|\eps|<1/4\). Apply \eqref{eq:udisplacement}.
\end{proof}

\begin{proposition}[Exact core duality and boundary/sector matching]
\label{prop:duality}
For all sufficiently large \(n\), the two core normalizations satisfy the
exact identity
\begin{equation}
 r_{\eps_n}\left(\frac1{2n}\right)=\frac1{B_n}.
 \label{eq:exact-duality}
\end{equation}
Moreover,
\begin{equation}
 \frac{n\rho_{\eps_n}^{\,n}u_n(\eps_n)}
 {U_{\eps_n,c}-U_{\eps_n}(\rho_{\eps_n}e^{-1/(2n)})}
 \longrightarrow\frac1{\sqrt{2\pi}}.
 \label{eq:boundary-sector}
\end{equation}
\end{proposition}
\begin{proof}
Substituting \(r=1/B_n\) into the core equation gives
\[
 \frac{c_n}{2B_n^2}h_{\eps_n}(B_n)=\frac1{2n},
\]
so uniqueness proves \eqref{eq:exact-duality}. The convergent chart gives
\(t(1/(2n))\sim1/B_n\). Also \(B_n/n\to0\), so the term \(\delta=1/(2n)\)
is negligible in \(t-\delta\). The numerator in
\eqref{eq:boundary-sector} equals \(\Pr(S_n=n)\sim1/(\sqrt{2\pi}B_n)\),
which proves the result.
\end{proof}

The exact identity \eqref{eq:exact-duality} explains why the same crossover
factor appears in local inversion and large-sector coefficients. Equation
\eqref{eq:boundary-sector} is a leading-order bridge, not an all-orders
transfer theorem. Constructing a uniform correction calculus on the Fourier
side is a distinct next step.

\section{Finite real bounds for inversion}
\label{sec:certificates}
The convergent chart has an existence-level geometric remainder. The
following finite inequalities offer an independent route to real enclosures.
They also avoid evaluating nearly cancelling gamma and zeta poles.

\begin{proposition}[Finite inversion inequalities]
\label{prop:finite}
For real \(|\eps|<1/4\), \(\alpha=2-\eps\), integer \(J\ge1\), and
\(t>0\), set
\begin{equation}
 H_J=\sum_{j=1}^Jj^{1-\alpha},\qquad
 K_J=\sum_{j=1}^Jj^{2-\alpha}.
 \label{eq:HK}
\end{equation}
Then
\begin{align}
 L_J(t):=c_\eps\left(\frac{t^2H_J}{2}-\frac{t^3K_J}{6}\right)
 &\le\Psi_\eps(t) \nonumber\\
 &\le c_\eps\left(\frac{t^2H_J}{2}
             +\frac{tJ^{1-\alpha}}{\alpha-1}\right)=:U_J(t),
 \label{eq:finite-bounds}\\
 c_\eps\left(tH_J-\frac{t^2K_J}{2}\right)
 &\le\Psi_\eps'(t)
 \le c_\eps\left(tH_J+\frac{J^{1-\alpha}}{\alpha-1}\right).
 \label{eq:finite-derivative}
\end{align}
If \(0<t_-<t_+\), and finite evaluations establish
\begin{equation}
 U_{J_-}(t_-)\le\delta\le L_{J_+}(t_+),
 \label{eq:root-certificate}
\end{equation}
then the root of \(\Psi_\eps(t)=\delta\) lies in \([t_-,t_+]\).
\end{proposition}
\begin{proof}
For all \(x\ge0\), Taylor's integral remainder gives
\[
 \frac{x^2}{2}-\frac{x^3}{6}\le e^{-x}-1+x\le\frac{x^2}{2},
 \qquad 0\le e^{-x}-1+x\le x,
\]
and
\(x-x^2/2\le1-e^{-x}\le x\), together with \(1-e^{-x}\le1\).
Apply the quadratic or cubic bounds to \(j\le J\), keep the omitted part
nonnegative for the lower bounds, and use
\[
 \sum_{j>J}j^{-\alpha}\le\int_J^\infty y^{-\alpha}\dd y
 =\frac{J^{1-\alpha}}{\alpha-1}
\]
for the upper tails. This proves both inequalities. Strict monotonicity of
\(\Psi_\eps\) proves the bracket assertion.
\end{proof}

If \(t_0\in[t_-,t_+]\) and a positive lower derivative bound \(m\) is
obtained at \(t_-\) from \eqref{eq:finite-derivative}, then convexity makes
\(m\) a lower derivative bound on the whole interval. Therefore
\begin{equation}
 |t-t_0|\le
 \frac{\max\{|L_J(t_0)-\delta|,|U_J(t_0)-\delta|\}}m.
 \label{eq:residual-certificate}
\end{equation}
Choosing \(J\) of order \(1/t\) is natural near the boundary; larger or
adaptively chosen finite sums can sharpen the bracket. These statements are
rigorous inequalities between real numbers. To turn them into numerical
certificates, the finite sums, powers, and \(c_\eps\) must themselves be
outward-rounded. The supplied numerical diagnostics do not do that, and are
not labeled interval certificates.

\section{Reproducible checks and numerical diagnostics}
\label{sec:checks}
The package contains \path{code/verify.py}, exact-check reports, CSV data, and
the generated tables below. The analytic proofs do not depend on the
numerical observations. The checks serve different purposes and are reported
separately rather than combined into a misleading ``verified theorem'' label.

\subsection{Exact finite algebra}
The program uses rational arithmetic to compare direct formal fixed-point
iteration with the Lagrange formula for
\[
 U(q)=\sum_{1\le j\le M}\frac{q^j e^{jU(q)}}{j^3}.
\]
This rational-amplitude test model deliberately omits the critical factor
\(c_0=1/\zeta(2)\). It checks the general algebraic identity, not a claim that
the critically normalized coefficients are rational. Degrees \(1\) through
\(12\), five cutoffs, positivity, and equality of coefficients below the
cutoff are checked exactly. A symbolic residual check independently verifies
\eqref{eq:P1}--\eqref{eq:P2}. The combined report records \textbf{122 passing
exact assertions}. These are finite computer-algebra checks, not a
proof-assistant formalization of the infinite statements.

\subsection{Positive coefficient recurrences}
For Poisson means \(\nu_j\), define \(p_k=\Pr(S=k)\). Differentiation of the
probability generating function gives the positive recurrence
\begin{equation}
 p_0=e^{-\sum_j\nu_j},\qquad
 kp_k=\sum_{j=1}^k j\nu_jp_{k-j}.
 \label{eq:panjer}
\end{equation}
The total intensity includes the full infinite tail even though only
\(j\le k\) enters the recurrence for \(p_k\). For a cutoff the intensity
and recurrence are both truncated consistently. The program evaluates this
recurrence through \(n=8192\), and forms the ratio using the exact expression
\eqref{eq:cutoff-identity}, with high-precision evaluations of the tail
intensity.

\begin{table}[htbp]
\centering\small
\begin{tabular}{rrrrrr}
\toprule
$n$ & $\lambda$ & $A_n/B_n$ & $\sqrt{2\pi}B_n\Pr(S_n=n)$ & $R_n(A_n)$ & $R_n(1.5A_n)$\\
\midrule
512 & -1 & 0.49813 & 0.83518 & 0.63957 & 0.87800 \\
512 & 0 & 0.53482 & 0.81565 & 0.61883 & 0.85704 \\
512 & 1 & 0.57793 & 0.76690 & 0.58099 & 0.83451 \\
2048 & -1 & 0.44669 & 0.86360 & 0.65661 & 0.86320 \\
2048 & 0 & 0.48270 & 0.84967 & 0.63216 & 0.85032 \\
2048 & 1 & 0.52348 & 0.81450 & 0.59662 & 0.83748 \\
8192 & -1 & 0.40890 & 0.88359 & 0.64122 & 0.85878 \\
8192 & 0 & 0.44420 & 0.87327 & 0.62845 & 0.84946 \\
8192 & 1 & 0.48330 & 0.84665 & 0.60452 & 0.83788 \\
\bottomrule
\end{tabular}

\caption{Coefficient diagnostics for \(\eps=\lambda/\log n\). Cutoff
arguments are rounded down to integers. The limiting values in the last
three columns are \(1\), \(e^{-1/2}\approx0.60653\), and
\(e^{-2/9}\approx0.80074\), respectively.}
\label{tab:coefficients}
\end{table}

The logarithmic convergence is slow, and the cutoff rounding is visible;
the table is not evidence for a monotone error law. Six separate comparisons
at \(n=128\), using an independent arbitrary-precision implementation of
\eqref{eq:panjer}, gave maximum relative disagreement
\(7.13\times10^{-18}\) or less. That tests the finite numerical calculation,
not its distance from the asymptotic limit.

The large-order recurrence uses NumPy's extended-range \texttt{longdouble}
when available. The program detects platforms with insufficient exponent
range and raises an error instead of silently underflowing. The exact checks
and the arbitrary-precision inverse checks do not require that extended
floating-point format. None of these calculations is outward-rounded.

\subsection{Direct checks of the convergent chart}
Choose \(r=e^{-L_r}\), \(\eps=\tau/L_r\), and set \(\delta\) by the exact
core. The program solves the original polylogarithm equation at 160 decimal
digits and compares \(t/r\) with \(1+P_1\) and \(1+P_1+P_2\).
All nine root residuals, measured as
\(|\Psi_\eps(t)/\delta-1|\), are below \(10^{-55}\).

\begin{table}[htbp]
\centering\small
\begin{tabular}{rrrrr}
\toprule
$\log(1/r)$ & $\varepsilon\log(1/r)$ & $t/r$ & degree-1 error & degree-2 error\\
\midrule
20 & -2 & 0.962316330 & $2.916\times10^{-3}$ & $2.837\times10^{-4}$ \\
20 & 0 & 0.963656348 & $1.156\times10^{-3}$ & $1.553\times10^{-5}$ \\
20 & 2 & 0.933247764 & $8.874\times10^{-3}$ & $2.230\times10^{-4}$ \\
40 & -2 & 0.980462618 & $7.624\times10^{-4}$ & $3.750\times10^{-5}$ \\
40 & 0 & 0.981541081 & $2.911\times10^{-4}$ & $1.888\times10^{-6}$ \\
40 & 2 & 0.968813278 & $2.248\times10^{-3}$ & $2.660\times10^{-5}$ \\
80 & -2 & 0.990045254 & $1.952\times10^{-4}$ & $4.824\times10^{-6}$ \\
80 & 0 & 0.990698010 & $7.301\times10^{-5}$ & $2.325\times10^{-7}$ \\
80 & 2 & 0.984965161 & $5.653\times10^{-4}$ & $3.246\times10^{-6}$ \\
\bottomrule
\end{tabular}

\caption{Direct inversion diagnostics. The error columns are absolute
errors in the ratio \(t/r\), not in \(t\). The two approximations are the
Taylor polynomials by total degree from \eqref{eq:P1}--\eqref{eq:P2}.}
\label{tab:inversion}
\end{table}

The observed improvement is consistent with the convergent chart but does not
supply a numerical value for the Cauchy constants in
\eqref{eq:chart-tail}. The reports distinguish a small evaluated residual
from an interval enclosure of the true root.

\subsection{Rebuilding and rerunning}
The source is self-contained apart from standard TeX packages and the two
included table fragments. From the package directory, run
\begin{verbatim}
python code/verify.py --part all
latexmk -pdf -interaction=nonstopmode Confluent_Critical_Transseries.tex
\end{verbatim}
The options \texttt{--part exact}, \texttt{--part coefficients}, and
\texttt{--part inversion} run the groups independently. Recorded package
versions, build instructions, and validation status are included alongside
the article. Python assertions should not be disabled with optimization flags
when running the check program.

\section{Further research questions and concrete targets}
\label{sec:future}

\paragraph{1. Uniform higher-order sector corrections.}
The local chart is convergent to all orders, whereas the triangular sector
law is leading-order. Expand the centered characteristic exponent beyond
\eqref{eq:K-compact} in a set of parameters that remains ordered when
\(\eps\log n\) crosses zero. The terms generated by
\(x^{2-\eps}\), \(x^2\), and their coalescing coefficients should be retained
in cancellation-free blocks. A useful target is a relative expansion uniform
for \(|\eps\log n|\le K\), with an explicit first correction and a remainder
small enough to distinguish it from the leading Gaussian term.

\paragraph{2. Quantitative cutoff errors and shrinking tolerances.}
Make the compact characteristic-function error and the two tail integrals
explicit in \(v_n,d_n\), and in
\(1-h_{\eps_n}(sA_n)/h_{\eps_n}(B_n)\). This would give a computable error
for \(R_n(M)\) on compact positive cutoff intervals. Then study
\(\theta_n\downarrow0\) in \eqref{eq:min-budget}. The fixed-tolerance
quantile argument does not itself justify a vanishing-tolerance budget.
An eventual interval implementation should certify a finite coefficient
ratio, not merely print its limiting profile.

\paragraph{3. The largest coupling window with untilted extremes.}
In Theorem~\ref{thm:coupling}, a bounded central displacement has no effect on
the limiting extreme process. If the target displacement is \(x_nB_n\),
a heuristic Gaussian density ratio for adding an action of size \(A_ny\)
contains \(\exp(x_nA_ny/B_n)\). This suggests that a new tilt can emerge
when \(|x_n|A_n/B_n\) is order one. Determine the actual validity range of
this heuristic, the necessary moderate-deviation estimates, and the resulting
conditional point process. The heuristic is not a theorem of this article.

\paragraph{4. General triangular slowly varying tails.}
Replace the exact tail by
\(c_nj^{-3+\eps_n}\ell_n(j)\). Identify verifiable uniform conditions under
which an implicit truncated-second-moment scale and a tail-quantile extreme
scale still satisfy the retained-variance property
\eqref{eq:truncated-variance}. Mere slow variation separately for each \(n\)
is unlikely to control the joint limit. The masked local-limit proof suggests
that this retained-variance property, an aperiodicity condition, and uniform
tail bounds are the correct inputs for a more general theorem.

\paragraph{5. Joint finite-fold drift.}
Truncating the kernel restores an analytic finite-action fold. Determine the
critical-point and critical-value drift uniformly in
\((\eps,M,n)\), particularly for \(M\asymp A_n\). Match the resulting
saddle formula to the cutoff factor \(H(M/A_n)\) with an error that remains
small after multiplication of logarithmic errors by \(n\). This would extend
the repository's fixed-exponent finite-fold analysis into the confluent window.

\paragraph{6. Complex continuation of the core.}
The analytic correction \(\cV\) is holomorphic in independent local
variables, but the parameterized core \(r_\eps(\delta)\) has its own complex
branch geometry. Construct explicit sectorial domains and track their
critical values uniformly as \(\eps\to0\). Only after doing so is it
reasonable to ask for uniform complex asymptotics, Stokes transitions, or
resurgent properties. A convergent correction chart alone does not establish
any of those global assertions.

\paragraph{7. Nonlinear slopes.}
For feedback of the form
\(U=\sum a_jq^j\exp((j+r_j)U)\), the substitution \(z=qe^U\) no longer
reduces the problem to a scalar polylogarithm. Start with
\(r_j=\gamma j^\theta\), \(0<\theta<1\), and determine whether it changes
the extreme scale, the Gaussian normalization, or only the critical centering.
A successful result must replace the exact coefficient/Poisson identity, not
use it after its hypotheses have been lost.

\paragraph{8. Signed weights and competing singularities.}
Positivity underlies cutoff monotonicity, the conditioning interpretation,
and several estimates. For signed or complex amplitudes, seek explicit
analytic hypotheses that permit a multi-saddle replacement. A particularly
concrete target is a finite signed prefix with a positive tail, where the
local normal form still exists but coefficient cancellation can invalidate
probability arguments.

\paragraph{9. Growing prefixes and the edge of robustness.}
Corollary~\ref{cor:prefix} keeps the altered prefix length fixed. Permit a
prefix length \(J_n\to\infty\), and characterize the allowable changes by
their centered second-moment contribution and their action-scale mass.
Perturbations at \(j\asymp A_n\) can change the limiting Poisson intensity
while leaving the central Gaussian limit unchanged; perturbations spread
through the variance-carrying range can change both. Prove a classification
that separates these effects quantitatively.

\paragraph{10. The lower endpoint \(\alpha\downarrow1\).}
The other half of the repository's endpoint question remains distinct. There
the first moment, not just the second, approaches divergence; the present
normalization \(c_\eps=1/\zeta(\alpha)\) tends to zero and the centering
mechanism changes. Derive an appropriate compensated core and determine
whether conditioning couples the largest actions to the total at leading
order. None of the upper-endpoint Gaussian estimates should be transported
there without a new proof.

\paragraph{11. Uniform matching to the fixed stable conditioned process.}
For fixed \(\alpha<2\), the natural fluctuation and largest-action scales
are comparable, while they separate in the joint limit proved here.
Construct an explicit normalization of the fixed stable bridge and show how
its conditioned jump process tends to the independent Poisson extreme
process as \(\alpha\uparrow2\). Then compare that limit with the present
arbitrary-rate theorem. This would explain the transition at the level of
entire conditioned processes, not only coefficient ratios.

\paragraph{12. Proof-producing computation and formalization.}
A practical first milestone is to formalize the finite identities
\eqref{eq:lagrange}, \eqref{eq:cutoff-identity}, and
\eqref{eq:finite-bounds}, while keeping their different algebraic and analytic
hypotheses explicit. Next build an interval evaluator for the real core and
finite root brackets. The harder Lean interfaces concern the masked Fourier
majorant, countable Poisson sums, and local conditioning. An implementation
should expose separately: an exact coefficient identity, a certified finite
numerical bound, and an asymptotic theorem with its hypotheses.

\section{Conclusion}
The upper endpoint is not obtained by inserting \(\alpha=2\) into a
fractional-power formula. The divergent gamma and zeta terms must first be
combined. Their exact combination is the elementary function
\(h_\eps\), which produces both a convergent inverse chart and the correct
triangular fluctuation scale.

For every rate of approach to exponent two, the extreme scale becomes small
relative to the fluctuation scale, yet a fixed multiple of the extreme scale
retains asymptotically all of the relevant truncated variance. This is why
exact conditioning preserves a Poisson extreme process and why the action
budget has a sharp explicit profile. The same normalization enters the local
inverse through the exact identity \(r_\eps(1/(2n))=1/B_n\).

These results resolve a specified upper-endpoint model problem left open in
the inspected research sequence. They supply new material for that sequence
without treating the fixed-endpoint classical mechanism as a new discovery,
without promoting numerical agreement to proof, and without conflating a
convergent local chart with global transseries realization or resurgence.

\appendix
\section{Source audit and dependency boundaries}
\label{app:provenance}

\subsection{Repository and library provenance}
The repository snapshot used for the final open-question crosswalk is
\begin{center}
\small\texttt{0c973d8f5e7ef4550f4df1bf486d890037fd9f14}.
\end{center}
The inspected project group is
\begin{center}
\small\path{Analysis/Transseries/docs/series-and-transseries/}.
\end{center}
The project README, group index, and the critical-Hahn article's research
questions and dependency discussion were read. The exact upper-endpoint
question was checked in the pinned article at source lines 1467--1472.
The full canonical transseries volume and every unmerged arrival were not
audited. Accordingly, this article makes a precise continuation claim for the
identified question, not an assertion that all possible repository overlap
has been excluded.

The user's Library supplied the scope statements of
\path{article(20260929-201315).tex}, titled \emph{The Logarithmic Critical
Endpoint}, and \path{Marginal_Critical_Transseries.tex}. These documents
identify the fixed endpoint and the still-unaddressed uniform crossover.
They are unpublished research drafts, not independent validation of the
present claims. No repository files were modified, and no Lean build of this
article's theorems was performed.

\subsection{Proof dependencies}
The analytic branch of the argument is
\[
 \text{polylogarithm expansion}
 \ \Longrightarrow\ \text{pole cancellation}
 \ \Longrightarrow\ \text{analytic implicit chart}.
\]
The coefficient branch is
\[
 \text{Lagrange identity}
 \ \Longrightarrow\ \text{Poisson representation}
 \ \Longrightarrow\ \text{masked local limit}.
\]
The scale estimates needed by the second branch follow directly from the
integral defining \(h_\eps\), not from an assumed stable-limit theorem.
Poisson tilting and tail Riemann sums then give the conditional process;
positivity and monotonicity give the sharp budget and its necessity.
The exact core duality connects the two branches after their independent
construction. The finite root bounds use only scalar exponential inequalities
and integral comparisons. All computational checks are downstream of these
proofs.

\section{A notation and regime ledger}
\small
\begin{longtable}{@{}p{0.23\textwidth}p{0.72\textwidth}@{}}
\toprule
Symbol & Meaning\\
\midrule
\endhead
\(\eps,\alpha\) & Moving parameter and exponent, related by \(\alpha=2-\eps\).\\[1pt]
\(c_\eps,F_\eps\) & Critical normalization \(1/\zeta(2-\eps)\) and normalized polylogarithm kernel.\\[1pt]
\(\rho_\eps,U_{\eps,c}\) & Critical radius \(e^{-F_\eps(1)}\) and boundary value \(F_\eps(1)\).\\[1pt]
\(\delta,t\) & Boundary coordinates \(q=\rho_\eps e^{-\delta}\), \(z=e^{-t}\).\\[1pt]
\(h_\eps,\cE\) & Exact crossover integral and entire divided exponential.\\[1pt]
\(r,v,\cV\) & Exact inverse core, its inverse-logarithmic parameter, and convergent analytic correction.\\[1pt]
\(A_n\) & Extreme-action scale \((nc_n)^{1/\alpha_n}\).\\[1pt]
\(B_n\) & Large root of \(B_n^2=nc_nh_{\eps_n}(B_n)\); central fluctuation scale.\\[1pt]
\(v_n,d_n\) & Small quantities \(h_{\eps_n}(B_n)^{-1}\) and \(nc_nB_n^{-\alpha_n}\).\\[1pt]
\(S_n,T_n\) & Full and thinned compound Poisson weighted sums.\\[1pt]
\(R_n(M),H(s)\) & Exact retained coefficient fraction and limiting profile \(e^{-1/(2s^2)}\).\\[1pt]
\(\Pi_n\) & Point process of action sizes divided by \(A_n\), with multiplicities.\\[1pt]
\(\lambda_n\) & Coefficient crossover parameter \(\eps_n\log n\).\\[1pt]
\(\omega_n\) & Bounded coupling displacement in \(\beta_n=1+\omega_nB_n/n\).\\
\bottomrule
\end{longtable}
\normalsize

\begin{center}
\small
\begin{tabular}{@{}lll@{}}
\toprule
Joint regime & Fluctuation normalization & Extreme normalization\\
\midrule
\(\eps\log n\to\lambda\in\R\) &
\(\sqrt{(c_0/2)n\log n\,\cE(\lambda/2)}\) &
\(\sqrt{nc_0}\,e^{\lambda/4}\)\\[5pt]
\(\eps\log n\to+\infty\) &
\((nc_\eps/\eps)^{1/(2-\eps)}\) &
\((nc_\eps)^{1/(2-\eps)}\)\\[5pt]
\(\eps\log n\to-\infty\) &
\(\sqrt{nc_\eps/|\eps|}\) &
\((nc_\eps)^{1/(2-\eps)}\)\\
\bottomrule
\end{tabular}
\end{center}
In every row \(\eps\to0\) is part of the hypothesis. The entries for \(B_n\)
are asymptotic equivalents; the general expression for \(A_n\) is exact.
The Gaussian local limit is not asserted for fixed positive \(\eps\).

\phantomsection
\addcontentsline{toc}{section}{References}
\begin{thebibliography}{9}

\bibitem{repo-index}
Vladimir Reshetnikov, \emph{ProveIt: Transseries project and
series-and-transseries index}, repository documentation, inspected revision
\texttt{0c973d8f5e7ef4550f4df1bf486d890037fd9f14}, 29 September 2026.
\url{https://github.com/VladimirReshetnikov/ProveIt/blob/0c973d8f5e7ef4550f4df1bf486d890037fd9f14/Analysis/Transseries/docs/series-and-transseries/README.md}.
Inspection scope and status caveats are described in Appendix~\ref{app:provenance}.

\bibitem{repo-stable}
\emph{Beyond Finite-Action Folds: Critical Hahn Transseries, Stable Sector
Asymptotics, and Sharp Action Budgets}, research package in \texttt{ProveIt},
29 September 2026. Research Question~4, source lines 1467--1472, inspected at
the revision above.
\url{https://github.com/VladimirReshetnikov/ProveIt/blob/0c973d8f5e7ef4550f4df1bf486d890037fd9f14/Analysis/Transseries/docs/series-and-transseries/Critical_Hahn_Transseries_Beyond_Finite_Action_Folds/article.tex}.

\bibitem{prior-endpoint}
\emph{The Logarithmic Critical Endpoint: Convergent Lambert Charts,
All-Order Sector Laws, and Sharp Action-Cutoff Corrections}, unpublished
research draft prepared for Vladimir Reshetnikov, 29 September 2026.
Inspected in the user's Library under
\path{article(20260929-201315).tex}. The scope statement explicitly treats
the fixed endpoint and excludes the joint exponent/order limit. Used for
scope comparison, not as a proof dependency.
% ed. (2026-09-29): filed location added.
Editorial note (ProveIt, 2026-09-29): now filed (batch~48) as
\path{Analysis/Transseries/docs/series-and-transseries/Logarithmic_Critical_Endpoint_Lambert_Charts/article.tex};
the joint limit treated here is its Question~1.

\bibitem{prior-marginal}
\emph{Marginal Critical Transseries: Convergent Lambert--W Charts,
Universal Cutoff Corrections, and Certified Action Budgets}, unpublished
research draft prepared for Vladimir Reshetnikov, 29 September 2026.
Inspected in the user's Library under
\path{Marginal_Critical_Transseries.tex}. Its conclusion identifies uniform
endpoint crossover as a further task. Used for scope comparison only.
% ed. (2026-09-29): filed location added.
Editorial note (ProveIt, 2026-09-29): now filed (batch~48) as
\path{Analysis/Transseries/docs/series-and-transseries/Marginal_Critical_Transseries_Action_Budgets/Marginal_Critical_Transseries.tex};
the joint limit treated here is its Question~1.

% ed. (2026-09-29): editorial bibliography entry for the editorial note in Section 1.1.
\bibitem{ed:sge}
[Editorial addition, ProveIt, 2026-09-29.]
\emph{Through the Stable--Gaussian Endpoint: Uniform Critical Transseries,
Prefix-Independent Cutoff Corrections, and Conditional Extremes}, research
package filed in batch~49:
\path{Analysis/Transseries/docs/series-and-transseries/Stable_Gaussian_Endpoint_Uniform_Critical_Transseries/article.tex}.

\bibitem{dlmf}
NIST Digital Library of Mathematical Functions, Chapter~25, especially
\S25.12(ii), Eq.~25.12.12, and the zeta functional equation in \S25.4.
\url{https://dlmf.nist.gov/25.12.E12}.
Accessed 29 September 2026. The standard polylogarithm expansion is the
external special-function identity used in Proposition~\ref{prop:normal}.

\bibitem{janson}
Svante Janson, Simply generated trees, conditioned Galton--Watson trees,
random allocations and condensation, \emph{Probability Surveys} \textbf{9}
(2012), 103--252. DOI: \texttt{10.1214/11-PS188}.
\url{https://arxiv.org/abs/1112.0510}.
Example~18.29 is a classical fixed cubic-tail antecedent of
Gaussian/extreme-value separation. The moving-exponent local-limit proof in
this article is given independently.

\end{thebibliography}
\end{document}
